% Section 3/11: evidence typology & protocol (P4)
Our claim standard replaces adjectives with ledger entries. Every reliability claim we make about
a method or system carries one of the following classes, applied at the level of individual
numbers, not whole papers.

\subsection{Evidence classes}
\begin{description}[leftmargin=1.5em, itemsep=1pt]
\item[[T]] Theoretical promise under explicit assumptions (conformal abstention, provable
alignment, identifiability results, distribution-free risk bounds). Such a claim carries its
assumption set forward, and the guarantee is conditional by construction.
\item[[E]] Independent empirical: evaluation of an artifact whose authors are not the evaluators,
whether the evaluators are a third party measuring what someone else built, shipped or not, or our
own re-verification of reported results against released artifacts.
\item[[V]] Vendor-reported: developer-documented claims from a model card, engineering note, or
commercial system description. Legitimate evidence of what a deployed system does; not evidence
that an unstudied analogous mechanism would transfer.
\item[[A]] Author-reported: results the artifact's authors report, whether or not they ran the
experiment (preprint tables, released evaluation-harness runs, self-evaluations of their own
prompt or method). Such a report documents what that team demonstrated; it is never independent
reproduction, and where artifacts are released we record that verification is checkable, not that it happened.
\item[[S]] Survey or synthesis where the aggregation is verified by us against the primary
records, or where our own citation-graph verification substitutes for reproduction. Our own
analytic cuts are [J], distinct from [S]: a placement or ordering we choose is judgement over the
record.
\item[[W]] Watch/pending: entities real enough to have acquired separate documentation while
our verification remains incomplete, or a specific result whose full-text check is not yet
complete (the card says which). We keep such entries explicit in the ledger so a reader can tell
what is solid from what is merely plausible.
\end{description}

\begin{table}[!t]
\centering
\small
\begin{tabular}{@{}l p{5.3cm} p{6.3cm}@{}}
\toprule
Token & One-line meaning & Test a reader applies \\
\midrule
\textsf{[T]} & theory under stated assumptions & does the claim carry its assumptions, and would the guarantee fail without them? \\
\textsf{[E]} & independent evaluation of someone else's artifact & did evaluators who did not author the artifact run the measurement? \\
\textsf{[A]} & results reported by the artifact's authors & is the measurement the authors' own report of their own artifact? \\
\textsf{[V]} & vendor or developer documentation & is the source the vendor's own description of a deployed system? \\
\textsf{[S]} & synthesis verified against the primary records & did we check the aggregation against those records? \\
\textsf{[W]} & watch/pending: verification incomplete & is the check still incomplete? \\
\textsf{[C]} & contradiction with a named anchor & is the contradicting anchor named? \\
\textsf{[N]} & named gap: a scoped, still-open absence & is the absence scoped to what was searched? \\
\textsf{[GAP]} & the specific unmeasured link behind a gap & which step has no anchor? \\
\textsf{[J]} & our analytic judgement; no chip & is this cut ours rather than a verified aggregation? \\
\bottomrule
\end{tabular}
\caption{The canonical evidence alphabet: six classes ([T], [E], [A], [V], [S], [W]) and four
annotations ([C], [N], [GAP], [J]). This is the single alphabet every other surface cites; no
surface defines a second one.}
\label{tab:evidence-classes}
\end{table}

Other surfaces list only the classes they use and defer to \cref{tab:evidence-classes}.

Separately from these classes, four annotations record what a card or caption is reacting to,
while the class carries how strong the evidence is.
[C] marks a specific contradiction (the anchor is named);
[N] is an absence flag --- a scoped, still-open gap in what has been measured, not a negative
finding; [GAP] marks the
specific unmeasured link behind such a gap, used where a card or a
table points at the step that has no anchor; and [J] marks our analytic judgement, used in figure
captions, figure annotations, and procedure steps where a cut is ours rather than a verified
aggregation. Where a claim carries two classes
we write them [T, E]: both apply, and the first listed is dominant. A bare tag placed directly
after a citation is also a valid citation-class badge: it classifies the cited result.

% Fig. 23 declared after the class list: at the section head the float cut the list between
% [S] and [W], so the last class resumed under the caption and read as a figure note.
\begin{figure}[!t]
  \centering
  % \includegraphics[width=\linewidth]{figures/f5_evidence_map.pdf}
  \resizebox{\linewidth}{!}{% figures/f5_evidencemap_draft.tex — F5 v4 (2026-10-19): V4 redesign D1 "ribbon ledger".
% One row per paper section (y), one rsRibSeg chip per claim family (x). Segment length therefore
% carries WHERE a family sits and how many the section holds; class is carried by the chip itself
% (fill weight + letter badge), so the old five-band x-grid -- and its 63mm index (41% of figure
% height) -- are gone. The 18 ids are the 16 first-pass families plus 17 (S4/L1, reasoning
% post-training) and 18 (S9, diffusion serving); both sit at the END of their row, so the rows
% still read in ledger order, and neither row grows past S8's five-chip row. The picture's natural
% width is set by the annotation lines below the ledger, not by any chip row, so nothing re-scales,
% no chip width changes, and the font floor is unchanged.
% 17 is [E] on the strongest anchor the family has (fuzzed checkers, FPR 0.83--0.87) rather than
% on the RLVR pass@1 gain, which is author-reported ([A]) and would have mislabelled the row's
% only [E]; 18 is [E] because the deployed-vs-declared gap rests on router telemetry, not on the
% vendor's declared rate.
% Contested is a STATE of a claim, not a class: ids 2, 11, 12 are the same [E] chip (green 16% fill)
% with a dashed vermillion BORDER -- the class is never replaced, the state rides the border only.
% [A] id 16 = purple fill, solid. Class key = one T-E-A-V-S row (five chips, same grammar), plus
% [C]/[N] annotations under the ledger. Thin rows (S5, S6) hold one family each: labelled, because
% the whitespace is a count, not a finding (D3).
% label: fig:evidence   (figure env + \label live in main.tex, sec:protocol)
%
% CAPTION NOTE — the main.tex caption still says "one hex = ..." and must be rewritten:
%   one chip = one claim family from the claim ledger; row = paper section, chip order inside a
%   row = ledger order (both are our placement [S]); chip length carries no quantity beyond the
%   count of families in that section; class = chip fill + letter badge ([T] outline, [E] filled
%   green, [A] filled pink, [V] outline pink, [S] filled slate); contested = the same chip,
%   dashed vermillion, for the three most contested families (2 JEPA-in-language stability,
%   11 monitor-evasion attackability, 12 calibration direction) -- surfaced, not hidden; delete
%   the "[W]" from that clause (the figure's state tag is [C], per H4).
%   No counts of papers, rates or sample sizes are encoded.
% Anchor companions (caption may list): #3 = 2605.26379 + 2606.12471 + 2607.22430; #4 = 2402.15991
%   + 2406.11353; #5 = 2305.13971 + 2501.10868 + 2411.15100; #1 = Laya repo TD card [V];
%   #14 = 2602.04101 [V]; #15 = CM test_3_verdict over 11 deep-reads [S]; #16 = 2610.14162 [A].
% Fonts: 7.2pt chips at natural ~157mm; published (resizebox to 165.1) ~7.6pt.
\begin{tikzpicture}[x=1mm,y=1mm]
\tikzset{
  rs/f5row/.style={anchor=east, align=right, font=\fontsize{7.5}{9.0}\selectfont, text=black!85},
  rs/f5note/.style={anchor=west, font=\fontsize{7.5}{9.0}\selectfont, text=black!70},
  rs/f5seg/.style={align=center, minimum width=24mm, minimum height=8.6mm,
    font=\fontsize{7.2}{8.4}\selectfont},
}
% ===== class key: one T-E-A-V-S row, same chip grammar as the ledger =====
\rsRibbonAt{0}{8.5}
\rsRibSeg{rsTheory}{\textbf{[T]} theory}
\rsRibSeg{rsVerified}{\textbf{[E]} empirical}
\rsRibSeg{rsAuthorRep}{\textbf{[A]} author}
\rsRibSeg{rsVendor}{\textbf{[V]} vendor}
\rsRibSeg{rsNoClaim}{\textbf{[S]} synthesis}
% ===== ledger rows (11mm pitch), sections top-down =====
% -- S4 L1 learning objectives: 1 [V], 2 [E] contested, 16 [A], 17 [E]
\rsRibbonAt{0}{0}
\rsRibSeg[rs/f5seg]{rsVendor}{\textbf{[V]} 1\\proper scoring}
\rsRibSeg[rs/f5seg, dashed, draw=rsContested]{rsVerified}{\textbf{[E]} 2\\JEPA stability}
\rsRibSeg[rs/f5seg]{rsAuthorRep}{\textbf{[A]} 16\\staged tokenizer}
\rsRibSeg[rs/f5seg]{rsVerified}{\textbf{[E]} 17\\verifier gaming}
% -- S5 L2 representations: 3 [T] only
\rsRibbonAt{0}{-11}
\rsRibSeg[rs/f5seg]{rsTheory}{\textbf{[T]} 3\\identifiability}
% -- S6 L3 inference architectures: 4 [E] only
\rsRibbonAt{0}{-22}
\rsRibSeg[rs/f5seg]{rsVerified}{\textbf{[E]} 4\\cascade gaps}
% -- S7 L4 output contracts: 5 [E], 6 [E], 7 [T]
\rsRibbonAt{0}{-33}
\rsRibSeg[rs/f5seg]{rsVerified}{\textbf{[E]} 5\\syntax checks}
\rsRibSeg[rs/f5seg]{rsVerified}{\textbf{[E]} 6\\typed decision}
\rsRibSeg[rs/f5seg]{rsTheory}{\textbf{[T]} 7\\risk selection}
% -- S8 L5 assurance: 8 [T], 9 [S], 10 [E], 11 [E] contested, 12 [E] contested
\rsRibbonAt{0}{-44}
\rsRibSeg[rs/f5seg]{rsTheory}{\textbf{[T]} 8\\conformal gate}
\rsRibSeg[rs/f5seg]{rsNoClaim}{\textbf{[S]} 9\\UQ inventory}
\rsRibSeg[rs/f5seg]{rsVerified}{\textbf{[E]} 10\\verifier fail}
\rsRibSeg[rs/f5seg, dashed, draw=rsContested]{rsVerified}{\textbf{[E]} 11\\monitor evasion}
\rsRibSeg[rs/f5seg, dashed, draw=rsContested]{rsVerified}{\textbf{[E]} 12\\calib. direction}
% -- S9 case studies: 13 [E], 14 [V], 18 [E]
\rsRibbonAt{0}{-55}
\rsRibSeg[rs/f5seg]{rsVerified}{\textbf{[E]} 13\\option flips}
\rsRibSeg[rs/f5seg]{rsVendor}{\textbf{[V]} 14\\hybrid gaps}
\rsRibSeg[rs/f5seg]{rsVerified}{\textbf{[E]} 18\\deployed vs declared}
% -- S11 review protocol: 15 [S]
\rsRibbonAt{0}{-66}
\rsRibSeg[rs/f5seg]{rsNoClaim}{\textbf{[S]} 15\\novelty verdict}
% ===== section labels (left) =====
\node[rs/f5row] at (-2,0)    {\S4 L1\\learning objectives};
\node[rs/f5row] at (-2,-11)  {\S5 L2\\representations};
\node[rs/f5row] at (-2,-22)  {\S6 L3\\inference architectures};
\node[rs/f5row] at (-2,-33)  {\S7 L4\\output contracts};
\node[rs/f5row] at (-2,-44)  {\S8 L5\\assurance};
\node[rs/f5row] at (-2,-55)  {\S9\\case studies};
\node[rs/f5row] at (-2,-66)  {\S11\\review protocol};
% ===== annotations: contested state, absent class, thin-row count =====
\node[rs/f5note] at (-2,-78)   {\textbf{[C]} contested: the same chip, dashed vermillion (ids 2, 11, 12).};
\node[rs/f5note] at (-2,-81.6) {\textbf{[N]} absence flag --- no family in this ledger is plotted as one.};
\node[rs/f5note] at (-2,-85.2) {thin rows (\S5, \S6) hold one family each; \S10 adds none --- counts, not findings.};
\end{tikzpicture}
}
  \caption{ Evidence ledger: one chip = one claim family from the claim ledger; the row is its paper section and the chip order inside a row is ledger order (both our analytic judgement [J]). Row length carries how many families that section contributes: one chip per family; no chip's width and no row's white space encodes a quantity of evidence. Class is carried by the chip: [T] outline, [E] filled green, [A] filled pink, [V] outline pink, [S] filled slate; chip class is the dominant class of the single claim it names, while a Table~1 row tag marks the strongest anchor of the family. The full alphabet, including [W] and the annotations, is \cref{tab:evidence-classes}. Contested stays an annotation: the same class chip, dashed vermillion border, for the three most contested families (2 JEPA-in-language stability, 11 monitor-evasion attackability, 12 calibration direction) --- the [C] annotation, surfaced rather than hidden. No family in the ledger is [N]; thin rows are counts, not findings. No rates, sample sizes or per-paper counts are encoded.}
  \label{fig:evidence}
\end{figure}

No evidence class admits a guarantee of universal correctness: [T] names assumption
conditions, [E] measures one deployment configuration, [V] reports one vendor's own design.

The classes encode provenance, not validity: a tag records who ran the measurement, never
whether the instrument measured the construct the citing sentence names. An [E] row can rest
on a contaminated benchmark, on a harness whose prompt format supplies half the effect, or on
a delta inside sampling noise, independently attributed and still not evidence for the claim
above it \cite{measwhatmatters2025,quantconstruct2026,contamsurvey2025,contamresist2026}.
Leaderboard position is the same exposure one level up, where rank rewards access to private
variants rather than reliability \cite{leaderillusion2025}. Where an anchor reports no
interval we print its number as the source printed it and do not read a point estimate as a
difference \cite{errorbars2024,evalci2026}.

\subsection{Claim-ledger construction}
For every claim that a card surfaces, the ledger records anchor (ID or artifact), verbatim-vs-paraphrase status, the numbers as written in their original context, the direction
(stats, benchmark structure) with the verification date, and the layer(s) the claim touches.
Any claim cited from sampling memory fails our protocol; every anchor is re-fetched live
(arXiv API verification with Microsoft/Elsevier content checks where applicable) before use in
this manuscript.

\subsection{Exclusions register}
We exclude claims from the ledger by a single rule: no claimworthy statement enters unless a
primary source stands behind it. The exclusions this cycle are concrete: a production rules-only / formally-verified pipeline (Case 6 is carried as an honest null on this very point); a JEPA-``SWM-distractors'' benchmark claim that failed verification and was never admitted to the ledger. The list continues with a structured-hallucination
benchmark without recoverable verification, a universal >20-option cost rule stated as a law (the vendor's per-system mechanism quote remains a [V] row in its own section),
the 0.85 ``default'' confidence threshold (replaced by threshold-as-policy), and a coinage for
which we could find no existing use. We treat it as ours to coin only if the definition appears
in our own terms. Each excluded row is retained verbatim in the build's review ledger so the
exclusion itself is auditable.

\subsection{Reproducibility record}
All searches ran the arXiv API (id-list/title/abs search, \(\le 10\) IDs per request) plus
publisher DOI metadata lookup, with snapshots per batch in the project's scan directory.
The paper's claim-to-evidence map is \cref{fig:evidence}, and every claim marked
contested in this manuscript gets its caveat in the same section.
Search cutoffs: 2026-10-17/18.
